\section{A mechanical CRC-routine check on the real 6-date parameter path}
\label{sec:crc-real-sequence}

Beyond the static, per-day comparison above, a narrower, purely
CONSTRUCTIVE question was checked directly on real numbers -- not a
statistical estimation, so it needs no forecasting power to be meaningful:
chaining the 6 real SPX dates' already-calibrated DLV$\times$Bates
backgrounds chronologically, does the REALISED change in $\theta$ and jump
parameters from one date to the next fall inside the reserve-CRC's
admissible ("permissible drift") region, and does \texttt{apply\_reset}
preserve protected-option prices exactly on these real numbers? Code:
\texttt{real\_data\_study/crc\_real\_sequence.py} (5/5 tests passing).

\paragraph{Within the standard region, 0 of 5 real transitions are
admissible.} Using the same conventions as the reserve-sweep study
($\pm0.03$ theta band, $1.5\times$-widened jump box): three transitions
fail because the realised $\theta$ move is 1.7--3.3$\times$ too large for
the declared band (all involve multi-year gaps); two fail because the
TARGET date's own jump-parameter calibration landed exactly on an
optimizer bound ($\lambda_J=5.0$ on 2024-03-05 and 2024-12-26;
$m_J=-0.4999$ on 2025-01-03) -- itself a flag that those specific fits are
poorly identified, not simply "large."

\paragraph{The mechanism is verified exactly correct on real numbers.} A
separate sensitivity pass widens each transition's region just enough to
admit the realised move (not a claim that this wider region is reasonable
-- a mechanism check only). Result: \texttt{apply\_reset} used the exact
realised $\Delta\theta$ with zero projection in every case, and protected-
option prices matched exactly (to floating-point zero) before vs.\ after
the reset on 4 of 5 transitions (Table~\ref{tab:crc-real-sequence}). The
5th hit a genuine numerical wall: admitting the bound-hitting $m_J=-0.4999$
forces a dominating envelope with $\delta_J\approx1.65$ (jump log-returns
with $\sim$165\% std), which overflows the Carr--Madan damped-contour
Fourier pricer at every quadrature setting tested -- a real fragility of
that pricing method for extremely wide envelopes, reported explicitly
rather than silently producing NaN.

\begin{table}[H]
\centering
\small
\begin{tabular}{llrrll}
\toprule
Transition & Gap & $\Delta\theta$ & $\theta$-widen$^\dagger$ & Blocked by (standard region) & Protected-price diff$^\ddagger$ \\
\midrule
2005-11-21 $\to$ 2010-07-07 & 4.62y & $+0.0988$ & 3.5x & $\theta$ move too large & $0.0$ (n=48) \\
2010-07-07 $\to$ 2011-07-07 & 1.00y & $-0.0519$ & 1.8x & $\theta$ move too large & $0.0$ (n=57) \\
2011-07-07 $\to$ 2024-03-05 & 12.66y & $+0.0274$ & 1.0x & jump params (target $\lambda_J$ at bound) & $0.0$ (n=51) \\
2024-03-05 $\to$ 2024-12-26 & 0.81y & $-0.0589$ & 2.1x & $\theta$ move too large & $0.0$ (n=18) \\
2024-12-26 $\to$ 2025-01-03 & 0.02y & $-0.0023$ & 0.1x & jump params (target $m_J$ at bound) & numerical failure$^\S$ (n=3) \\
\bottomrule
\end{tabular}

\vspace{2mm}
\footnotesize
$^\dagger$ Factor by which the $\pm0.03$ theta band would need to widen for admissibility (forced-admissible sensitivity pass; jump-box widening required was $1.0$x the standard $1.5$x factor in every case, i.e.\ no extra widening beyond the construction's own target-containment margin). $^\ddagger$ Maximum absolute price difference on protected contracts (expiry $<$ gap) before vs.\ after \texttt{apply\_reset}, once the region was forced wide enough to admit the realised move -- $0.0$ confirms exact price preservation on real numbers. $^\S$ Carr-Madan Fourier pricing overflowed for every quadrature setting tested; a genuine numerical limit of the dominating-envelope construction when forced to contain an extreme (bound-hitting) target jump mean, not resolved by tuning.

\caption{The 5 real chronological transitions: realised move, admissibility
under the standard region, and the forced-admissible mechanism check.}
\label{tab:crc-real-sequence}
\end{table}

\subsection{More data is needed for this to become a genuine historical-dynamics result}
\label{sec:more-data-needed}

This check is deliberately mechanical, not statistical, and that is both
its strength and its limit. It required no forecasting power to be
informative -- but it also cannot, by itself, establish whether a
\emph{properly fitted} physical law and reset-proposal policy would track
real SPX dynamics well, because:

\begin{itemize}
\item \textbf{The theta-band failures look like a time-spacing artefact,
  fixable by denser data.} The shortest real gap available (2024-12-26 to
  2025-01-03, $\sim$1 week) produced by far the smallest realised
  $\Delta\theta$ ($-0.0023$, comfortably inside the band on its own),
  consistent with $\theta$ moving roughly in proportion to elapsed time.
  Weekly or monthly real data would very plausibly bring most transitions
  inside a modestly-sized band.
\item \textbf{The jump-box failures are NOT primarily a spacing issue.}
  All three bound-hitting calibrations (2024-03-05, 2024-12-26,
  2025-01-03) are exactly the dense, modern dates where
  Sec.~\ref{sec:main-comparison} already found DLV$\times$Bates fit quality
  collapsing (13--27\% inside bid-ask) -- the same underlying model-
  capacity problem surfacing twice. Sampling MORE modern dates would
  likely surface more of the same degenerate fits, not fewer. Fixing this
  needs model/calibration changes (looser or better-justified bounds, more
  model flexibility, or a joint/warm-started sequential fit across the
  path instead of independent day-by-day re-calibration), not denser
  sampling alone.
\end{itemize}

\textbf{This environment has only 6 real, non-adjacent option-quote dates}
(2005--2025, gaps from 7 days to 13 years; see
DATA\_PROVENANCE.md). That is enough for the cross-sectional comparison
above and for the mechanical check in this section, but not enough to fit
or validate a genuine historical transition law, physical risk premia, or
an OCE/hedging challenger test with any statistical power -- at most one
real training transition and one held-out forecast point exist in the
closest-spaced real subsequence. A full historical-dynamics estimation
(latent-state filtering, an explicit $P$/$Q$ relation with
$b^P-b^Q=\Sigma\lambda$, admissible $\theta$-reset proposal dynamics fit
jointly across the path, chronological train/validation/held-out testing,
and held-out OCE/hedging tests) is specified and ready to run as soon as a
denser real series is available -- see
\texttt{real\_data\_study/HISTORICAL\_SIMULATOR\_READY\_PROMPT.md} for the
exact task specification to resume verbatim, and
\texttt{DATA\_PROVENANCE.md}/\texttt{DATA\_STATUS.md} for exactly what
data access (WRDS credentials or a denser pre-populated
\texttt{QUANTDATA\_ROOT}) would unblock it.
